\subsubsection*{Minimal Complete Bases}

% BEGIN GENERATED ARTIFACT: def:minimal-complete-basis-propositional-logic
\begin{definitionbox}{Definition (Minimal Complete Basis)}
\begin{definition}[Minimal Complete Basis]
\label{def:minimal-complete-basis-propositional-logic}
A functionally complete connective basis \(B\) is \emph{minimal complete} if no
proper subset of \(B\) is functionally complete.
\end{definition}
\end{definitionbox}

\begin{remark*}[Standard quantified statement]
\[
\operatorname{MinimalComplete}(B)
\Longleftrightarrow
\operatorname{FunctionallyComplete}(B)
\land
(\forall C\subsetneq B)\neg\operatorname{FunctionallyComplete}(C).
\]
\end{remark*}

\begin{remark*}[Predicate reading]
\(\operatorname{MinimalComplete}(B)\) means that \(B\) is complete, but removing
any primitive connective from \(B\) destroys completeness.
\end{remark*}

\begin{remark*}[Interpretation]
Minimality is relative to deletion of connectives from the basis. It does not
mean that formulas using the basis are short or easy to read.
\end{remark*}

\begin{dependencies}
\begin{itemize}
  \item \hyperref[def:functionally-complete-connective-set-propositional-logic]{Functionally Complete Connective Set}
  \item \hyperref[def:connective-basis-propositional-logic]{Connective Basis}
\end{itemize}
\end{dependencies}
% END GENERATED ARTIFACT: def:minimal-complete-basis-propositional-logic

% BEGIN GENERATED ARTIFACT: thm:nand-nor-minimal-complete-bases-propositional-logic
\begin{theorembox}{Theorem (NAND and NOR are Minimal Complete Bases)}
\begin{theorem}[NAND and NOR are Minimal Complete Bases]
\label{thm:nand-nor-minimal-complete-bases-propositional-logic}
The one-connective bases \(\{\uparrow\}\) and \(\{\downarrow\}\) are minimal
complete bases.

\noindent\hyperref[prf:nand-nor-minimal-complete-bases-propositional-logic]{Go to proof.}
\end{theorem}
\end{theorembox}

\begin{remark*}[Standard quantified statement]
\[
\operatorname{MinimalComplete}(\{\uparrow\})
\land
\operatorname{MinimalComplete}(\{\downarrow\}).
\]
\end{remark*}

\begin{remark*}[Predicate reading]
NAND alone is complete, NOR alone is complete, and deleting the sole connective
from either basis leaves no functionally complete basis.
\end{remark*}

\begin{remark*}[Interpretation]
A one-connective complete basis is automatically minimal once the empty basis is
not functionally complete. NAND and NOR are therefore minimal complete bases
because each is complete by itself.
\end{remark*}

\begin{dependencies}
\begin{itemize}
  \item \hyperref[def:minimal-complete-basis-propositional-logic]{Minimal Complete Basis}
  \item \hyperref[thm:nand-functionally-complete-propositional-logic]{NAND is Functionally Complete}
  \item \hyperref[thm:nor-functionally-complete-propositional-logic]{NOR is Functionally Complete}
  \item \hyperref[def:incomplete-connective-set-propositional-logic]{Incomplete Connective Set}
\end{itemize}
\end{dependencies}
% END GENERATED ARTIFACT: thm:nand-nor-minimal-complete-bases-propositional-logic
